\documentclass[10pt,aspectratio=169,serif]{beamer}

\usepackage[fontset=none, scheme=plain]{ctex}
\IfFontExistsTF{Songti SC}{%
	\setCJKmainfont{Songti SC}[AutoFakeSlant=0.2]%
	\setCJKsansfont{Songti SC}[AutoFakeSlant=0.2]%
}{%
	\setCJKmainfont{FandolSong-Regular}[Extension=.otf,BoldFont=FandolSong-Bold,AutoFakeSlant=0.2]%
	\setCJKsansfont{FandolSong-Regular}[Extension=.otf,BoldFont=FandolSong-Bold]%
}
\usepackage{amsmath,amssymb,graphicx,tikz,booktabs,multicol,calligra}
\usepackage{pgfplots}
\pgfplotsset{compat=1.18}
\usepackage{tcolorbox}
\tcbuselibrary{skins}
\usepackage{XyygU}
\usepackage{caption}
\captionsetup[figure]{name=图,labelsep=period}

\usetikzlibrary{arrows.meta,positioning,calc}

% 公式间距
\setlength{\abovedisplayskip}{0pt plus 2pt minus 2pt}
\setlength{\belowdisplayskip}{0pt plus 2pt minus 2pt}

\title[集合间的基本关系练习]{集合间的基本关系（练习）}
\subtitle{子集、真子集、集合相等}
\author{}
\institute{}
\date{2026-10-16}

\begin{document}

\begin{frame}
  \titlepage
\end{frame}

% ========== 第1题 ==========
\begin{frame}[t]{第1题：写出所有子集}
\textbf{1.} 写出集合 $\{a,\ b,\ c\}$ 的所有子集。
\end{frame}

\begin{frame}[t]{第1题 解答}
\begin{tcolorbox}[colback=orange!10,colframe=orange,title=\textbf{解}]
含有 $n$ 个元素的集合共有 $2^n$ 个子集。

集合 $\{a, b, c\}$ 含 $3$ 个元素，故共有 $2^3 = 8$ 个子集，按元素个数列举如下：

\vspace{0.5em}
\begin{itemize}
  \item 空集（$0$ 个元素）：$\varnothing$；
  \item 单元素集（$1$ 个元素）：$\{a\}$，$\{b\}$，$\{c\}$；
  \item 双元素集（$2$ 个元素）：$\{a,b\}$，$\{a,c\}$，$\{b,c\}$；
  \item 三元素集（$3$ 个元素）：$\{a,b,c\}$。
\end{itemize}

\vspace{0.5em}
\textbf{综上：} $\varnothing$，$\{a\}$，$\{b\}$，$\{c\}$，$\{a,b\}$，$\{a,c\}$，$\{b,c\}$，$\{a,b,c\}$。

\vspace{0.3em}
{\small 其中除 $\{a,b,c\}$ 本身外，其余 $7$ 个均为它的真子集。}
\end{tcolorbox}
\end{frame}

% ========== 第2题 ==========
\begin{frame}[t]{第2题：用适当符号填空}
\textbf{2.} 用适当的符号填空（$\in$，$\notin$，$=$，$\subseteq$，$\subsetneqq$ 等）：

\vspace{0.8em}
\begin{columns}
\begin{column}{0.55\textwidth}
(1) $a$ \underline{\hspace{1em}} $\{a,\ b,\ c\}$；

(3) $\varnothing$ \underline{\hspace{1em}} $\{x \in \mathbf{R} \mid x^2 + 1 = 0\}$；

(5) $\{0\}$ \underline{\hspace{1em}} $\{x \mid x^2 = x\}$；
\end{column}
\begin{column}{0.45\textwidth}
(2) $0$ \underline{\hspace{1em}} $\{x \mid x^2 = 0\}$；

(4) $\{0,\ 1\}$ \underline{\hspace{1em}} $\mathbf{N}$；

(6) $\{2,\ 1\}$ \underline{\hspace{1em}} $\{x \mid x^2 - 3x + 2 = 0\}$。
\end{column}
\end{columns}
\end{frame}

\begin{frame}[t]{第2题 解答}
\begin{tcolorbox}[colback=orange!10,colframe=orange,title=\textbf{解}]
\textbf{(1)} $a \in \{a,b,c\}$（$a$ 是集合中的元素，用属于号）。

\vspace{0.3em}
\textbf{(2)} $0 \in \{x \mid x^2 = 0\}$（方程 $x^2 = 0$ 的解为 $x=0$）。

\vspace{0.3em}
\textbf{(3)} $\varnothing = \{x \in \mathbf{R} \mid x^2 + 1 = 0\}$（方程 $x^2+1=0$ 无实根，该集合为空集）。

\vspace{0.3em}
\textbf{(4)} $\{0, 1\} \subsetneqq \mathbf{N}$（$0,1$ 都是自然数，但自然数还有 $2,3,\dots$）。

\vspace{0.3em}
\textbf{(5)} $\{0\} \subsetneqq \{x \mid x^2 = x\}$（$x^2=x$ 的解为 $x=0$ 或 $x=1$，右边为 $\{0,1\}$）。

\vspace{0.3em}
\textbf{(6)} $\{2, 1\} = \{x \mid x^2 - 3x + 2 = 0\}$（方程的解为 $x=1$ 或 $x=2$，两边元素相同）。
\end{tcolorbox}
\end{frame}

% ========== 第3题 ==========
\begin{frame}[t]{第3题：判断集合间的关系}
\textbf{3.} 判断下列两个集合之间的关系：

\vspace{0.5em}
(1) $A = \{x \mid x < 0\}$，$B = \{x \mid x < 1\}$；

\vspace{0.5em}
(2) $A = \{x \mid x = 3k,\ k \in \mathbf{N}\}$，$B = \{x \mid x = 6z,\ z \in \mathbf{N}\}$；

\vspace{0.5em}
(3) $A = \{x \in \mathbf{N}_+ \mid x \text{ 是 } 4 \text{ 与 } 10 \text{ 的公倍数}\}$，
$B = \{x \mid x = 20m,\ m \in \mathbf{N}_+\}$。
\end{frame}

\begin{frame}[t]{第3题 解答}
\begin{tcolorbox}[colback=orange!10,colframe=orange,title=\textbf{解}]
\textbf{(1)} 若 $x < 0$，则必有 $x < 1$，即 $A$ 中元素都在 $B$ 中；但 $0.5 \in B$ 而 $0.5 \notin A$。

\alert{故 $A \subsetneqq B$。}

\vspace{0.5em}
\textbf{(2)} $A$ 是 $3$ 的倍数（$0,3,6,9,12,\dots$），$B$ 是 $6$ 的倍数（$0,6,12,18,\dots$）。$6$ 的倍数一定是 $3$ 的倍数，反之不成立（如 $3 \in A$ 但 $3 \notin B$）。

\alert{故 $B \subsetneqq A$。}

\vspace{0.5em}
\textbf{(3)} $4$ 与 $10$ 的最小公倍数为 $20$，故 $A = \{20, 40, 60, \dots\}$；
而 $B = \{20m \mid m \in \mathbf{N}_+\} = \{20, 40, 60, \dots\}$。两集合元素完全相同。

\alert{故 $A = B$。}
\end{tcolorbox}
\end{frame}

\begin{frame}[c]
  \begin{center}
    {\Huge\calligra Thanks!}
  \end{center}
\end{frame}

\end{document}
